\section{Introduction}
\label{sec:intro}

Colonoscopy is the gold standard for colorectal cancer screening.
Computer-aided detection (CAD) systems analyze individual frames to flag suspected polyps and to confirm anatomical landmarks (cecum, ileum) required for quality metrics.
Deep convolutional networks are attractive for near real-time deployment, yet they require large labeled datasets and often overfit a single hospital's appearance.

Self-supervised learning can reduce annotation cost by pretraining on unlabeled video frames.
Joint-Embedding Predictive Architectures (JEPA)~\cite{lecun2022path,assran2023ijepa} learn by predicting latent representations of masked views instead of reconstructing pixels.
Medical JEPA variants exist for ultrasound and MRI, but not for colonoscopy with lightweight CNN backbones suitable for clinical deployment.

\enbrief{We provide an open, reproducible pipeline: unsupervised CNN-JEPA pretraining on colonoscopy frames, mask-aware SE blocks compatible with JEPA masking, and a three-class fine-tuning benchmark (landmark / normal mucosa / polyp) with label-efficiency and leave-one-center-out protocols.
We do \emph{not} claim state-of-the-art accuracy over ImageNet on public Simula data; we report results honestly and release code.}

\textbf{Contributions.}
(1)~\textbf{Method:} a convolutional JEPA (\method{}) on \backbone{} for unlabeled endoscopy frames.
(2)~\textbf{Technique:} mask-aware SE blocks using mask-weighted global average pooling under spatial masking.
(3)~\textbf{Evaluation:} reproducible C0--C10 benchmark on three clinical classes with early stopping, label fractions (10/25/50\%), and LOO framework (currently one public center).
